% ==================================================
% MODEL II: SET-BASED LOCAL OPTIMIZATION
% ==================================================

\section{Model II: Set-Based Local Optimization}
\label{sec:set}

Now, we need a second model that will use the valid timetable produced by the OR-tree, and produce an output minimizing the soft penalties applied due to the soft constraints we have in the problem.

Set-based search works in this case because it allows for small improvements on the solution pass by pass.

% ==================================================
\subsection{Search model: $A_{set}=(S_{set},T_{set})$}
% ==================================================

Our search model is $A_{set}=(S_{set},T_{set})$. We first establish the facts that make up our timetables and the rules that allow us to modify them.

\paragraph{Facts.}

Following Lecture 06, we define our set of facts as:

\begin{equation}
    F=\{\Assigned(x,s)\mid x\in X,\ s\in D(x)\}
\end{equation}

Each fact is simply one assignment in our timetable. By combining these facts, we can represent an entire timetable as a single set.

\paragraph{States.}

Lecture 06 defines the state space as a subset of the power set of facts:

\begin{equation}
    S_{set}\subseteq 2^F
\end{equation}

For our problem, we only want to consider states that represent complete, valid timetables. Therefore, we define:

\begin{equation}
    S_{set}=
    \left\{
    q\subseteq F\ \middle|\
    \begin{array}{l}
    q\text{ assigns every }x\in X\text{ exactly once,}\\
    \mathrm{Constr}(a_q)=\mathrm{true}
    \end{array}
    \right\}
\end{equation}

Here, $q$ represents a set of assignment facts, while $a_q$ is the actual timetable represented by those facts. Every state in our search is already a valid solution; at this stage we improve assignments rather than schedule more components.

\paragraph{Extensions.}

Following Lecture 06, extensions replace facts: $\EXT\subseteq\{A\longrightarrow B\mid A,B\subseteq F\}$. We propose two types of extensions: Move and Swap.

\textbf{1. Move}

A Move extension changes the time slot assigned to a single component.

\begin{equation}
    \{\Assigned(x,s_1)\}
    \longrightarrow
    \{\Assigned(x,s_2)\}
\end{equation}

Here, $x$ is the component being moved, $s_1$ is its current time slot, and $s_2$ is the new time slot.

For example, we could move CPSC 231 LEC 02 from MWF 9--10 to MWF 10--11, provided that doing so does not violate any hard constraints.

\textbf{2. Swap}

A Swap extension exchanges the assigned slots of two components.

\begin{equation}
    \begin{aligned}
    &\{\Assigned(x,s_1),\Assigned(y,s_2)\}\\
    &\qquad\longrightarrow
    \{\Assigned(x,s_2),\Assigned(y,s_1)\}
    \end{aligned}
\end{equation}

Here, $x$ and $y$ are two distinct components, while $s_1$ and $s_2$ are their current slots.

We use swaps because sometimes moving one component on its own may not be possible, but exchanging the slots of two components could still produce a better timetable.

For both extensions, we need to ensure that the new assignments belong to the components' allowed domains, respect preassignments, and do not violate any hard constraints.

This is important because we do not want our optimization process to undo the work already done by the OR-tree.

\paragraph{Transitions.}

Once we select an extension rule, we can apply it to the current state using the transition defined in Lecture 06:

\begin{equation}
    q'=(q\setminus A)\cup B
\end{equation}

Our transition relation is therefore:

\begin{equation}
    \begin{aligned}
    T_{set}=\{(q,q')\mid\;&
    \exists(A\rightarrow B)\in\EXT,\\
    &A\subseteq q,\\
    &q'=(q\setminus A)\cup B,\\
    &q'\in S_{set}\}
    \end{aligned}
\end{equation}

The last condition, $q'\in S_{set}$, ensures that every transition leads to another valid timetable.

% ==================================================
\subsection{Search process: $P_{set}=(A_{set},\Env_{set},K_{set})$}
% ==================================================

Following Lecture 06, our search process is $P_{set}=(A_{set},\Env_{set},K_{set})$.

\paragraph{Environment.}

Our environment contains the allowed slot domains, hard constraints, preferences and penalty values; its current configuration records $q$ along with that fixed instance data.

\paragraph{Search control.}

For our control, we propose a greedy hill-climbing approach.

The idea is fairly simple. At every step, we consider the available Move and Swap extensions, calculate the resulting penalties, and choose the extension that gives us the lowest evaluation score.

However, we only apply an extension if it actually improves the current timetable.

\textbf{Weight function}

Lecture 06 uses $f_{weight}$ to assign a numerical value to each applicable extension rule.

For our problem, we define:

\begin{equation}
    f_{weight}(A,B,e)=\Eval\left(a_{(q(e)\setminus A)\cup B}\right),
\end{equation}
where $q(e)$ is the current state recorded in environment $e$.

The function essentially asks: if we perform this Move or Swap, what will the new timetable's penalty be?

Since our objective is to minimize penalties, a lower weight is preferable.

\textbf{Applicable extensions}

We define the set of currently applicable extensions as:

\begin{equation}
    \begin{aligned}
    R(q)=\{(A,B)\mid\;&
    A\rightarrow B\in\EXT,\\
    &A\subseteq q,\\
    &(q\setminus A)\cup B\in S_{set},\\
    &(q\setminus A)\cup B\neq q\}
    \end{aligned}
\end{equation}

For nonempty $R(q)$, let
\[
w_{\min}(q)=\min_{(A,B)\in R(q)}f_{weight}(A,B,e).
\]

\textbf{Selection function}

It is possible that multiple extensions produce the same minimum weight.

In that case, we use the function $f_{select}$ from Lecture 06 to break the tie.

First, we collect the best extensions:

\begin{equation}
    R_{\min}(q)=
    \{(A,B)\in R(q)\mid
    f_{weight}(A,B,e)=w_{\min}(q)\}
\end{equation}

We then select one:

\begin{equation}
    (A^*,B^*)=
    f_{select}(R_{\min}(q),e)
\end{equation}

Our proposed tie-breaker simply chooses randomly among the equally weighted extensions.

This gives us one extension to consider applying.

\textbf{Strict improvement}

We only apply an extension when $w_{\min}(q)<\Eval(a_q)$. Then
\[
K_{set}(q,e)=(q\setminus A^*)\cup B^*.
\]
Otherwise, we stop the search and return the current timetable.

\paragraph{Local optimum.}
A local optimum is a valid timetable where no single Move or Swap can reduce the penalty. This does not necessarily mean we have found the best timetable possible.

% ==================================================
\subsection{Search instance: $\Inst_{set}=(q_0,G_{set})$}
% ==================================================

The search instance is $\Inst_{set}=(q_0,G_{set})$.

\paragraph{Initial state.}

Since the OR-tree has already produced a valid timetable, we do not need to start with an empty state.

Let $a_{\mathrm{feasible}}$ represent the timetable returned by the OR-tree.

We convert it into a set of assignment facts:

\begin{equation}
    q_0=
    \{\Assigned(x,a_{\mathrm{feasible}}(x))
    \mid x\in X\}
\end{equation}

Since $\mathrm{Constr}(a_{\mathrm{feasible}})=\mathrm{true}$, we know $q_0\in S_{set}$.

\paragraph{Goal condition.}

Our goal is to stop at zero penalty, or when no improving neighbor exists. Let
\[
N(q)=\{q'\in S_{set}:(q,q')\in T_{set},\ q'\neq q\}.
\]
Then
\[
G_{set}(q)=\yes\iff \Eval(a_q)=0\ \lor\
\neg\exists q'\in N(q):\Eval(a_{q'})<\Eval(a_q).
\]
For larger instances, we could also introduce a time or iteration limit.

% ==================================================
\subsection{Implementation and operation}
We start with the OR-tree's valid timetable, generate all possible valid neighbors through Move and Swap extensions, and select the best improvement. An implementation stores $q$ as a \textbf{set of assignment facts} $\Assigned(x,s)$ (equivalently, a dictionary from component to slot), $D(x)$ as a \textbf{dictionary of permitted slot lists}, and current room occupancies in \textbf{arrays or dictionaries}. Candidate Moves and Swaps can be enumerated using the component and slot lists, checked with $\mathrm{Consistent}$ and scored using $\Eval$. Only the current state, the best candidate and its score are required; the search need not retain a tree of past states. If no candidate reduces the penalty, the current timetable is returned.

\subsection{Worked set-based derivation}
To demonstrate how our search works, consider a simplified scheduling problem involving two lecture sections and two tutorials for CPSC 231.
\[
\begin{gathered}
L_1=\text{CPSC 231 LEC 01},\quad L_2=\text{CPSC 231 LEC 02},\\
T_1=\text{CPSC 231 TUT 01},\quad T_2=\text{CPSC 231 TUT 02}.
\end{gathered}
\]
Let $s_9,s_{10},s_{11}$ denote MWF 9--10, 10--11, 11--12;
$r_{10}$ denote TT 10--11, and $w_{11}$ denote MW 11--12. The domains are
\[
\begin{aligned}
D(L_1)&=\{s_9\}, & D(L_2)&=\{s_9,s_{10},s_{11}\},\\
D(T_1)&=\{r_{10}\}, & D(T_2)&=\{w_{11}\}.
\end{aligned}
\]
We assume lecture room capacity two, tutorial room capacity one, and both tutorials shared by both lecture sections. The only nonzero soft penalty is 5 if both lecture sections occupy the same slot:
\[
\Eval(a)=5\cdot\mathbf{1}[a(L_1)=a(L_2)].
\]
Suppose OR-tree search returns
\[
\begin{aligned}
q_0=\{&\Assigned(L_1,s_9),\Assigned(L_2,s_9),\\
&\Assigned(T_1,r_{10}),\Assigned(T_2,w_{11})\},
\qquad \Eval(a_{q_0})=5.
\end{aligned}
\]
Both lectures can share the slot under the assumed room capacity, and neither tutorial overlaps them.

The controller considers the Move extension
\[
\underbrace{\{\Assigned(L_2,s_9)\}}_{A}\longrightarrow
\underbrace{\{\Assigned(L_2,s_{10})\}}_{B}.
\]
This extension is valid because the new slot belongs to $D(L_2)$ and does not overlap either tutorial. Moving $L_2$ to $s_{11}$ is not permitted, since that slot overlaps $T_2$ on Monday and Wednesday. Applying the extension yields
\[
q_1=(q_0\setminus A)\cup B,
\qquad \mathrm{Constr}(a_{q_1})=\mathrm{true},
\qquad \Eval(a_{q_1})=0<5.
\]
Therefore our hill-climbing control accepts the extension and stops at zero penalty:
\[
q_0\xrightarrow{\text{Move }L_2:s_9\to s_{10}}q_1,
\qquad G_{set}(q_1)=\yes.
\]
This example shows how the set-based search takes a valid timetable from the OR-tree, applies a small change, and produces a better solution without violating any hard constraints.

% ==================================================
\subsection{Justification and limitations}
We chose set-based search because it fits naturally with the second stage of our scheduling problem. Rather than rebuilding the timetable, we can focus on improving the assignments through smaller changes.

Our Move and Swap extensions also give us a straightforward way to explore different schedules, while the weight function allows us to compare them using the evaluation function defined in the problem.

However, the main limitation comes from our greedy search control.

Hill climbing only accepts extensions that improve the current timetable. This means we may eventually reach a local optimum, even when a better timetable exists elsewhere.

There is also the possibility that certain improvements require multiple components to move together. Since our current extensions only move one component or swap two, the search may not always be able to reach those solutions.

We could address this later through more advanced controls, such as simulated annealing, or by introducing extension rules that modify several assignments at once.

For our proposed model, however, greedy hill climbing gives us a relatively simple optimization process that follows the set-based search framework from Lecture 06 while keeping every intermediate timetable valid.
